% Source: https://artofproblemsolving.com/community/c3h1534292_factor_this
% Stable identity: aops-topic-1534292-q01
\textbf{No integer factorisation.} If $p$ were a product of two nonconstant polynomials with integer coefficients, one of them would be linear, $x-c$ up to sign, and $c$ would be an integer root dividing the constant term $-7$. But
\[
p(1)=-1,\qquad p(-1)=-9,\qquad p(7)=455,\qquad p(-7)=-273,
\]
so there is no integer root, and $p$ cannot be factored over the integers. (By the rational root theorem it has no rational root either, so it is irreducible over the rationals.)

\textbf{Exactly one real root.} For real numbers $a<b$,
\[
p(b)-p(a)=(b^3-a^3)+2(b^2-a^2)+3(b-a)=(b-a)\left(a^2+ab+b^2+2a+2b+3\right),
\]
and completing the square twice,
\[
a^2+ab+b^2+2a+2b+3=\left(a+\dfrac {b}{2}+1\right)^2+\dfrac{3}{4}\left(b+\dfrac{2}{3}\right)^2+\dfrac{5}{3}>0 .
\]
Hence $p$ is strictly increasing, so it has at most one real root. Since $p(1)=-1<0<p(2)=15$ and $p$ is continuous, there is exactly one real root $r$, with $1<r<2$ (numerically $r\approx1.10$; in the sketch the vertical scale is compressed by a factor of $4$):

\begin{asy}
usepackage("amsmath");
size(5cm);
defaultpen(fontsize(10pt));
import graph;
real f(real x){return (x^3+2x^2+3x-7)/4;}
draw((-2.2,0)--(2.4,0),Arrow(TeXHead));draw((0,-3.3)--(0,3),Arrow(TeXHead));
draw(graph(f,-1.8,1.75),blue+linewidth(1));
label("$x$",(2.4,0),E);label("$y$",(0,3),N);
real r0=1.0954; dot((r0,0),red); label("$r$",(r0,0),2.2*dir(-30),red);
label("$1$",(1,0),NW);draw((1,-0.07)--(1,0.07));label("$2$",(2,0),S);draw((2,-0.07)--(2,0.07));
label("$-7$",(0,-1.75),NW);draw((-0.06,-1.75)--(0.06,-1.75));
label("$y=p(x)$",(1.75,f(1.75)),E,blue);
label("\mbox{}",(min(currentpicture,user=true)+max(currentpicture,user=true))/2,black); // keeps the text after the figure black when LaTeX typesets the labels inline
\end{asy}

So the equation has one real root and two non-real complex roots, and the polynomial does not factor over the integers. $\blacksquare$
